\documentclass[10pt,a4paper]{amsart}
\usepackage{amsmath,amssymb,amsthm,graphicx}
% Retained preamble for the Stacks Project chapter {CH} (shared project preamble).
% The chapter file pinned-original-source.tex only carries \input{{preamble}} on its first line,
% so no preamble text lives inside the chapter itself.  The Stacks Project uses one single shared
% preamble for every chapter: the snapshot stages it three times (sources/stacks-chow/preamble.tex,
% sources/stacks-etale-cohomology/preamble.tex and focused-seeds/AGP/sources/stacks-snapshot-20261011/
% preamble.tex) and those copies are identical in size.  The block below is copied mechanically,
% line for line, from sources/stacks-chow/preamble.tex lines 14-18 and 154-217: the Xy-pic setup and
% the native theorem declarations together with the mathematical macros.  Nothing was added, removed
% or redefined; the wrapper class and the amsmath/amssymb/amsthm/graphicx base packages are supplied
% by the assembler.
% For commutative diagrams we use Xy-pic
% Chapter: stacks-spaces-more-morphisms.tex -- More on Morphisms of Spaces.
\usepackage[all]{xy}

% We use 2cell for 2-commutative diagrams.
\xyoption{2cell}

% Theorem environments.
%
\theoremstyle{plain}
\newtheorem{theorem}[subsection]{Theorem}
\newtheorem{proposition}[subsection]{Proposition}
\newtheorem{lemma}[subsection]{Lemma}

\theoremstyle{definition}
\newtheorem{definition}[subsection]{Definition}
\newtheorem{example}[subsection]{Example}
\newtheorem{exercise}[subsection]{Exercise}
\newtheorem{situation}[subsection]{Situation}

\theoremstyle{remark}
\newtheorem{remark}[subsection]{Remark}
\newtheorem{remarks}[subsection]{Remarks}

\numberwithin{equation}{subsection}

% Macros
%
\def\lim{\mathop{\mathrm{lim}}\nolimits}
\def\colim{\mathop{\mathrm{colim}}\nolimits}
\def\Spec{\mathop{\mathrm{Spec}}}
\def\Hom{\mathop{\mathrm{Hom}}\nolimits}
\def\Ext{\mathop{\mathrm{Ext}}\nolimits}
\def\SheafHom{\mathop{\mathcal{H}\!\mathit{om}}\nolimits}
\def\SheafExt{\mathop{\mathcal{E}\!\mathit{xt}}\nolimits}
\def\Sch{\mathit{Sch}}
\def\Mor{\mathop{\mathrm{Mor}}\nolimits}
\def\Ob{\mathop{\mathrm{Ob}}\nolimits}
\def\Sh{\mathop{\mathit{Sh}}\nolimits}
\def\NL{\mathop{N\!L}\nolimits}
\def\CH{\mathop{\mathrm{CH}}\nolimits}
\def\proetale{{pro\text{-}\acute{e}tale}}
\def\etale{{\acute{e}tale}}
\def\QCoh{\mathit{QCoh}}
\def\Ker{\mathop{\mathrm{Ker}}}
\def\Im{\mathop{\mathrm{Im}}}
\def\Coker{\mathop{\mathrm{Coker}}}
\def\Coim{\mathop{\mathrm{Coim}}}

% Boxtimes
%
\DeclareMathSymbol{\boxtimes}{\mathbin}{AMSa}{"02}

%
% Macros for moduli stacks/spaces
%
\def\QCohstack{\mathcal{QC}\!\mathit{oh}}
\def\Cohstack{\mathcal{C}\!\mathit{oh}}
\def\Spacesstack{\mathcal{S}\!\mathit{paces}}
\def\Quotfunctor{\mathrm{Quot}}
\def\Hilbfunctor{\mathrm{Hilb}}
\def\Curvesstack{\mathcal{C}\!\mathit{urves}}
\def\Polarizedstack{\mathcal{P}\!\mathit{olarized}}
\def\Complexesstack{\mathcal{C}\!\mathit{omplexes}}
% \Pic is the operator that assigns to X its picard group, usage \Pic(X)
% \Picardstack_{X/B} denotes the Picard stack of X over B
% \Picardfunctor_{X/B} denotes the Picard functor of X over B
\def\Pic{\mathop{\mathrm{Pic}}\nolimits}
\def\Picardstack{\mathcal{P}\!\mathit{ic}}
\def\Picardfunctor{\mathrm{Pic}}
\def\Deformationcategory{\mathcal{D}\!\mathit{ef}}
\usepackage{cleveref}
\crefname{originaldefinition}{original definition}{original definitions}
\Crefname{originaldefinition}{Original definition}{Original definitions}
\makeatletter
\newlabel{definition-geometrically-reduced-fibre}{{0E08}{source}{}{sourcepointer-1}{}}
\newlabel{definition-geometrically-reduced-fibre@cref}{{[originaldefinition][0][]0E08}{[1][0][]source}{}{}{}}
\makeatother
\crefname{originallemma}{original lemma}{original lemmas}
\Crefname{originallemma}{Original lemma}{Original lemmas}
\makeatletter
\newlabel{lemma-proper-flat-geom-red}{{0E1E}{source}{}{sourcepointer-2}{}}
\newlabel{lemma-proper-flat-geom-red@cref}{{[originallemma][0][]0E1E}{[1][0][]source}{}{}{}}
\makeatother
\crefname{originallemma}{original lemma}{original lemmas}
\Crefname{originallemma}{Original lemma}{Original lemmas}
\makeatletter
\newlabel{lemma-stein-universally-closed}{{0A19}{source}{}{sourcepointer-3}{}}
\newlabel{lemma-stein-universally-closed@cref}{{[originallemma][0][]0A19}{[1][0][]source}{}{}{}}
\makeatother
\crefname{originallemma}{original lemma}{original lemmas}
\Crefname{originallemma}{Original lemma}{Original lemmas}
\makeatletter
\newlabel{lemma-stein-universally-closed-residue-fields}{{0E1C}{source}{}{sourcepointer-4}{}}
\newlabel{lemma-stein-universally-closed-residue-fields@cref}{{[originallemma][0][]0E1C}{[1][0][]source}{}{}{}}
\makeatother
\crefname{originalreference}{original reference}{original references}
\Crefname{originalreference}{Original reference}{Original references}
\makeatletter
\newlabel{more-morphisms-lemma-etale-makes-integral-split}{{0BSR}{source}{}{sourcepointer-5}{}}
\newlabel{more-morphisms-lemma-etale-makes-integral-split@cref}{{[originalreference][0][]0BSR}{[1][0][]source}{}{}{}}
\makeatother
\crefname{originalreference}{original reference}{original references}
\Crefname{originalreference}{Original reference}{Original references}
\makeatletter
\newlabel{spaces-perfect-lemma-proper-flat-geom-red-connected}{{0E0S}{source}{}{sourcepointer-6}{}}
\newlabel{spaces-perfect-lemma-proper-flat-geom-red-connected@cref}{{[originalreference][0][]0E0S}{[1][0][]source}{}{}{}}
\makeatother
\crefname{originaltheorem}{original theorem}{original theorems}
\Crefname{originaltheorem}{Original theorem}{Original theorems}
\makeatletter
\newlabel{theorem-stein-factorization-general}{{0A1C}{source}{}{sourcepointer-7}{}}
\newlabel{theorem-stein-factorization-general@cref}{{[originaltheorem][0][]0A1C}{[1][0][]source}{}{}{}}
\makeatother
\title{The Stein factorization of a proper flat morphism with geometrically reduced fibres is finite etale}
\author{The Stacks Project Authors}
\date{stacks-project:spaces-more-morphisms.tex:local-snapshot-20261011; GNU FDL 1.2 or later}
\begin{document}
\maketitle
\noindent\textit{Editorial note.} Counted span is the author's own proof environment, source lines 8436-8476, immediately following the statement at 8425-8434.\par\medskip
\noindent\textit{Editorial note.} The statement cites the Stein factorization theorem and the definition of geometrically reduced fibres by \textbackslash{}ref; both are navigation pointers.\par\medskip
\noindent\textit{Editorial note.} Two cross-chapter pointers (More on Morphisms, Derived Categories of Spaces) are recorded in assembly.external\_references.\par\medskip
\noindent\textit{Allowed standard prerequisites (editorial).}\begin{itemize}
\item Stein factorization in general. chapter theorem theorem-stein-factorization-general, cited in the statement. Supplies the factorisation X -> Y' -> Y whose second map is shown to be finite etale.
\item Compatibility of the Stein factorization with flat base change. chapter lemma lemma-stein-universally-closed, cited at proof line 8438. Allows the fibre of Y' -> Y to be computed after a flat base change.
\item Proper flat morphisms with geometrically reduced fibres have geometrically connected fibres nearby. chapter lemma lemma-proper-flat-geom-red, cited at proof lines 8468-8469. Upgrades connectivity of one geometric fibre to connectivity of the fibres in a neighbourhood.
\item The residue field version of the Stein argument. chapter lemma lemma-stein-universally-closed-residue-fields, cited in the proof. Shows the fibre of Y' -> Y is a single point with the residue field of the base point.
\item Structure sheaf of a proper flat morphism with geometrically connected fibres. Derived Categories of Spaces, lemma spaces-perfect-lemma-proper-flat-geom-red-connected, cited at proof lines 8473-8474. Gives f\_*O\_X = O\_Y, which identifies Y' with Y locally and yields etaleness.
\end{itemize}\medskip
\ifdefined\hypertarget\hypertarget{sourcepointer-1}{}\fi
\noindent\textit{Source reference.} Original definition 0E08: sources/stacks-spaces-more-morphisms/pinned-original-source.tex, source lines 6904--6904. Navigation pointer only: the counted proof cites this statement by \textbackslash{}ref as a standing input; the deduction the counted proof has to make is carried out inside the counted span, and no part of that deduction is delegated to this pointer.\par\medskip
\ifdefined\hypertarget\hypertarget{sourcepointer-2}{}\fi
\noindent\textit{Source reference.} Original lemma 0E1E: sources/stacks-spaces-more-morphisms/pinned-original-source.tex, source lines 8387--8387. Navigation pointer only: the counted proof cites this statement by \textbackslash{}ref as a standing input; the deduction the counted proof has to make is carried out inside the counted span, and no part of that deduction is delegated to this pointer.\par\medskip
\ifdefined\hypertarget\hypertarget{sourcepointer-3}{}\fi
\noindent\textit{Source reference.} Original lemma 0A19: sources/stacks-spaces-more-morphisms/pinned-original-source.tex, source lines 7994--7994. Navigation pointer only: the counted proof cites this statement by \textbackslash{}ref as a standing input; the deduction the counted proof has to make is carried out inside the counted span, and no part of that deduction is delegated to this pointer.\par\medskip
\ifdefined\hypertarget\hypertarget{sourcepointer-4}{}\fi
\noindent\textit{Source reference.} Original lemma 0E1C: sources/stacks-spaces-more-morphisms/pinned-original-source.tex, source lines 8061--8061. Navigation pointer only: the counted proof cites this statement by \textbackslash{}ref as a standing input; the deduction the counted proof has to make is carried out inside the counted span, and no part of that deduction is delegated to this pointer.\par\medskip
\ifdefined\hypertarget\hypertarget{sourcepointer-5}{}\fi
\noindent\textit{Source reference.} Original reference 0BSR: focused-seeds/geometry/S9-spaces-stein-factorization-etale/ref-tags.txt, source lines 4--4. Navigation pointer only: the counted proof cites this statement by \textbackslash{}ref as a standing input; the deduction the counted proof has to make is carried out inside the counted span, and no part of that deduction is delegated to this pointer.\par\medskip
\ifdefined\hypertarget\hypertarget{sourcepointer-6}{}\fi
\noindent\textit{Source reference.} Original reference 0E0S: focused-seeds/geometry/S9-spaces-stein-factorization-etale/ref-tags.txt, source lines 5--5. Navigation pointer only: the counted proof cites this statement by \textbackslash{}ref as a standing input; the deduction the counted proof has to make is carried out inside the counted span, and no part of that deduction is delegated to this pointer.\par\medskip
\ifdefined\hypertarget\hypertarget{sourcepointer-7}{}\fi
\noindent\textit{Source reference.} Original theorem 0A1C: sources/stacks-spaces-more-morphisms/pinned-original-source.tex, source lines 8197--8197. Navigation pointer only: the counted proof cites this statement by \textbackslash{}ref as a standing input; the deduction the counted proof has to make is carried out inside the counted span, and no part of that deduction is delegated to this pointer.\par\medskip
\begin{lemma}
\label{lemma-stein-factorization-etale}
Let $S$ be a scheme.
Let $f : X \to Y$ be a proper morphism of algebraic spaces over $S$.
Let $X \to Y' \to Y$ be the Stein factorization of $f$
(Theorem \ref{theorem-stein-factorization-general}).
If $f$ is of finite presentation, flat, with geometrically
reduced fibres (Definition \ref{definition-geometrically-reduced-fibre}),
then $Y' \to Y$ is finite \'etale.
\end{lemma}

\begin{proof}
Formation of the Stein factorization commutes with flat base change,
see Lemma \ref{lemma-stein-universally-closed}.
Thus we may work \'etale locally on $Y$ and we may assume $Y$
is an affine scheme. Then $Y'$ is an affine scheme and $Y' \to Y$
is integral.

\medskip\noindent
Let $y \in Y$. Set $n$ be the number of connected components of
the geometric fibre $X_{\overline{y}}$. Note that $n < \infty$
as the geometric fibre of $X \to Y$ at $y$ is a proper
algebraic space over a field, hence Noetherian,
hence has a finite number of connected components.
By Lemma \ref{lemma-stein-universally-closed-residue-fields}
there are finitely many points $y'_1, \ldots, y'_m \in Y'$ lying over $y$
and for each $i$ we can pick a finite type point $x_i \in |X_y|$
mapping to $y'_i$ the extension $\kappa(y'_i)/\kappa(y)$ is finite.
Thus More on Morphisms,
Lemma \ref{more-morphisms-lemma-etale-makes-integral-split}
tells us that after replacing $Y$ by an \'etale neighbourhood
of $y$ we may assume $Y' = V_1 \amalg \ldots \amalg V_m$ as a scheme
with $y'_i \in V_i$ and $\kappa(y'_i)/\kappa(y)$ purely inseparable.
In this case the algebraic spaces $X_{y_i'}$
are geometrically connected over $\kappa(y)$, hence $m = n$.
The algebraic spaces $X_i = (f')^{-1}(V_i)$, $i = 1, \ldots, n$
are proper, flat, of finite presentation, with geometrically
reduced fibres over $Y$. It suffices to prove the lemma
for each of the morphisms $X_i \to Y$. This reduces us to the case where
$X_{\overline{y}}$ is connected.

\medskip\noindent
Assume that $X_{\overline{y}}$ is connected. By
Lemma \ref{lemma-proper-flat-geom-red}
we see that $X \to Y$ has geometrically connected
fibres in a neighbourhood of $y$. Thus
we may assume the fibres of $X \to Y$ are geometrically connected.
Then $f_*\mathcal{O}_X = \mathcal{O}_Y$ by
Derived Categories of Spaces, Lemma
\ref{spaces-perfect-lemma-proper-flat-geom-red-connected}
which finishes the proof.
\end{proof}

\end{document}
